% ============================================================
%  SECCIÓN 1: INTRODUCCIÓN DE LA EMPRESA Y ORGANIGRAMA
%
%  RECORDATORIO AL EQUIPO:
%  - Descripción breve pero completa: sector, tamaño, historia reciente,
%    presencia geográfica, número de empleados, cifras clave.
%  - El organigrama debe ser claro: incluir niveles jerárquicos, áreas
%    funcionales y, si es posible, el departamento de RRHH destacado.
%  - Captad cómo está estructurada la empresa formalmente vs. cómo
%    funciona realmente (cultura, informalidad, matricial vs. funcional).
%  - Si tenéis el organigrama como imagen, ponedla en /imagenes/
%    y usad \includegraphics. Si lo hacéis en LaTeX, usad tikz o forest.
% ============================================================

\section{Introducción de la empresa y organigrama}

% ---- 1.1 Descripción general ----
\subsection{Descripción general de CaixaBank}

% TODO: Ampliar con datos obtenidos de la entrevista y fuentes oficiales.
% Fuentes útiles: Memoria Anual CaixaBank, CNMV, web corporativa.

CaixaBank, S.A. es una entidad bancaria española con sede central en Valencia.
Surgida de la reorganización del Grupo ``la Caixa'', cotiza en el IBEX~35 y es
actualmente uno de los mayores grupos financieros de España y Portugal tras la
fusión con Bankia en 2021.

\begin{table}[H]
  \centering
  \caption{Datos corporativos básicos de CaixaBank (datos orientativos; actualizad con los reales)}
  \begin{tabularx}{\textwidth}{lX}
    \toprule
    \textbf{Parámetro}           & \textbf{Valor} \\
    \midrule
    Sector                       & Servicios financieros y bancarios \\
    Año de fundación             & 1904 (origen en Caja de Pensiones para la Vejez y de Ahorros) \\
    Sede central                 & Valencia, España \\
    Número de empleados          & $\sim$45.000 (grupo consolidado, 2024) \\
    Número de oficinas           & $\sim$4.000 en España \\
    Presencia internacional      & Portugal (BPI), Luxemburgo, Reino Unido, entre otros \\
    Facturación / Margen bruto   & [Dato real a incluir] \\
    \bottomrule
  \end{tabularx}
\end{table}

% ---- 1.2 Organigrama ----
\subsection{Organigrama}

% OPCIÓN A — Imagen externa (recomendada si tenéis el organigrama oficial):
% Guardad el archivo como "organigrama.png" (o .pdf) en /imagenes/ y
% descomentad las líneas siguientes:
%
% \begin{figure}[H]
%   \centering
%   \includegraphics[width=0.95\textwidth]{imagenes/organigrama}
%   \caption{Organigrama de CaixaBank (fuente: [indicar fuente])}
%   \label{fig:organigrama}
% \end{figure}

% OPCIÓN B — Descripción textual (mientras no tengáis la imagen):

A continuación se describe la estructura organizativa de CaixaBank de manera
resumida, pendiente de completar con el organigrama gráfico oficial:

\begin{itemize}
  \item \textbf{Consejo de Administración}: máximo órgano de gobierno corporativo.
  \item \textbf{Comisión Ejecutiva}: presidida por el Presidente Ejecutivo; supervisa
        la gestión diaria.
  \item \textbf{Dirección General}: organizada en grandes áreas funcionales:
    \begin{itemize}
      \item Área Financiera
      \item Área Comercial y Marketing
      \item Área de Tecnología y Datos
      \item \textbf{Área de Personas y Organización (RRHH)}
      \item Área de Riesgos
      \item Área Jurídica y Cumplimiento Normativo
    \end{itemize}
  \item \textbf{Red Territorial}: direcciones territoriales que agrupan las
        oficinas por zonas geográficas.
\end{itemize}

\begin{notaequipo}
  Completad este apartado con el organigrama real obtenido en la entrevista o
  en la web/memoria corporativa. Indicad si la estructura es \textbf{funcional},
  \textbf{divisional} o \textbf{matricial}, y si ha cambiado tras la fusión con
  Bankia. Preguntad también si existe un departamento de RRHH centralizado o
  si está descentralizado por territorios.
\end{notaequipo}
